\section{结论或计算心得}

本文在 Beans 三分类任务上，对 ResNeXt50\_32x4d 和 DenseNet121 的从头训练与微调方案进行了比较。综合 Accuracy、Macro-F1 与训练时间来看，微调在本实验设置下整体更有优势，尤其是在训练初期表现出更快的收敛速度。这说明在小样本图像分类任务中，ImageNet 预训练权重所提供的通用特征仍然具有较强的迁移价值。

从训练过程中的实际体会来看，从头训练并不是不能取得较好的结果，但它对超参数更敏感，训练前期也更容易出现震荡。如果只看最终结果，可能会觉得两种方法差别不算特别大；但真正跑完整个实验后，会发现微调不仅更容易达到较高指标，而且调参压力也相对更小。尤其是在 CPU 环境下，重复试验的成本比较高，这时稳定性的重要性会更加突出。

同时，本实验也说明模型结构的差异会影响训练表现。DenseNet 依靠特征复用，在小规模数据集上往往更容易表现出稳定性；ResNeXt 的结构表达能力较强，但在训练成本和收敛节奏上也可能更依赖具体设置。对课程实验而言，这一点给我的启发是：选择模型时不能只看“结构是否更先进”，还需要结合数据规模、硬件条件以及调参成本一起考虑。

本次实验仍有一定局限。例如，数据集规模较小，且每组实验只使用单次随机种子结果，没有继续做更充分的重复实验；另外，实验环境为 CPU，训练时间相对较长，也限制了更大规模的消融尝试。后续如果条件允许，可以增加更多数据集或加入更多 CNN 结构，进一步分析不同任务中迁移学习收益的变化情况。
